\documentclass{article}
\usepackage[margin=1in]{geometry}
\usepackage{amsmath, amssymb, amsthm}
\usepackage{enumitem}
\usepackage{mathtools}
\usepackage{cancel}

% colored links
\usepackage{hyperref}
\hypersetup{
    colorlinks=true,
    linkcolor=blue,
    filecolor=magenta,      
    urlcolor=black!20!BurntOrange,
    }

% Inputting Python code
\usepackage[dvipsnames]{xcolor}
\definecolor{textblue}{rgb}{.2,.2,.7}
\definecolor{textred}{rgb}{0.54,0,0}
\definecolor{textgreen}{rgb}{0,0.43,0}
\usepackage{upquote}
\usepackage{listings}
\lstset{
    language=Python, 
    tabsize=4,
    basicstyle={\ttfamily},
    keywordstyle=\color{textblue},
    commentstyle=\color{textgreen},
    stringstyle=\color{textred},
    frame=none,
    columns=fullflexible,
    keepspaces=true,
    showstringspaces=false,
    xleftmargin=-15mm, % manual adjustment, figure out permanent solution
}

% Drawing figures
\usepackage{tikz}
\usetikzlibrary{calc, shapes.symbols}

% Colored Boxes
\usepackage{tcolorbox}
\tcbuselibrary{skins,hooks}
\usetikzlibrary{shadows}
\usepackage{lipsum}

%Images
\usepackage{graphicx}
    \usepackage{subcaption}
    \usepackage{float}

%Formatting and Spacing
\setitemize[1]{noitemsep, parsep = 5pt, topsep = 5pt}
\setenumerate[1]{label = (\alph*), parsep = 1pt, topsep = 5pt}
\setlength\parindent{0pt}
\linespread{1.1}

%Easy Numbering in case we add or remove questions
\newcounter{counter}
\newcommand{\Exercise}{Exercise \stepcounter{counter}\thecounter}

\newcommand{\Cov}{\text{Cov}}
\newcommand{\trace}{\text{trace}}
\newcommand{\Normal}{\text{Normal}}
\newcommand{\var}[1]{\operatorname{var}\left(#1\right)}
\newcommand{\E}[1]{\operatorname{E}\left[#1\right]}
\renewcommand{\Pr}[1]{\operatorname{P}\left(#1\right)}

\DeclareMathOperator*{\minimize}{minimize}
\DeclareMathOperator*{\maximize}{maximize}
\DeclareMathOperator*{\argmin}{argmin}
\DeclareMathOperator*{\argmax}{argmax}
\DeclarePairedDelimiter\abs{\lvert}{\rvert}%
\DeclarePairedDelimiter\norm{\lVert}{\rVert}%
\DeclarePairedDelimiterX{\inp}[2]{\langle}{\rangle}{#1, #2}

%Custom Title Fields
\newcommand{\hmwkTitle}{Homework 4}
\newcommand{\hmwkDueDate}{Feb.\ 14, 2024}
\newcommand{\hmwkDueTime}{11:59pm EST}
\newcommand{\hmwkClass}{Advanced Algorithms}
\newcommand{\hmwkClassInstructor}{Professor Dana Randall}
\newcommand{\hmwkSection}{Spring 2024}
\newcommand{\hmwkAuthorName}{}

%Headers and Footers
\usepackage{fancyhdr}
\usepackage{extramarks}
\pagestyle{fancy} 
\lhead{CS 4540}
\chead{\hmwkClass \ (\hmwkClassInstructor)}
\rhead{\hmwkTitle}
\cfoot{\thepage}
\renewcommand\headrulewidth{0.4pt}
\renewcommand\footrulewidth{0.4pt}

\title{
    \vspace{2in}
    \textbf{\hmwkClass:\ \hmwkTitle}\\
    \normalsize\vspace{0.1in}\small{Due\ on\ \hmwkDueDate\ at \hmwkDueTime}\\
    \vspace{0.1in}\large{\textit{\hmwkClassInstructor\ \hmwkSection}} \\
    \vspace{0.8in}
    \begin{minipage}[t]{0.4\columnwidth}
        {\footnotesize\textbf{As stated in the syllabus, unauthorized use of previous semester course materials is strictly prohibited in this course.
        }}
    \vspace{0.5in}
    \end{minipage}
    \vspace{0.8in}
    \author{\textbf{\hmwkAuthorName}}
    \date{}
}

\begin{document}

\maketitle
\pagebreak

\textbf{Note: This should have already been inferred, but unless otherwise stated, justification for a result / design should always be given.}

\subsection*{\Exercise}
  For some problem $P$, two online algorithms $A_1$ and $A_2$ are given, with a
  competitive ratio of 2 and 3, respectively. Design a randomized online
  algorithm with competitive ratio $9/4$.

\subsection*{\Exercise}
  We consider a version of the game Memory with just one player. As in the
  usual game, let $n$ pairs of cards lie on the table, face down. In each move
  of the game, the player can turn two cards, one at a time. If they are
  identical, they are removed, otherwise they are turned face down again and
  put back in their spots. The goal of the player is to remove all cards with
  the least number of moves. Even after the successful turning of a pair, any
  further turning is considered a move. We assume that the player can remember
  all cards that have been turned at any time. The optimal strategy
  needs $n$ moves to remove all $2n$ cards.
  \begin{enumerate}
    \item Design an online strategy that takes at most $2n$ moves. i.e. that is 2-competitive
    \item Design an online strategy that takes at most $2n-1$ moves
    \item Show that any deterministic strategy can be forced to take at least $2n-1$ moves.
    \item What changes about the competitive ratio if the player gets a free move after removing a pair?
  \end{enumerate}

\subsection*{\Exercise}
  For the Paging Problem, let \textsc{Alg} be any marking algorithm as
  presented in the lecture with a cache of size k, and let \textsc{Opt} be an optimal
  off-line-algorithm with a cache of size $h \leq k$. Prove that \textsc{Alg} is
  $\frac{k}{k-h+1}$-competitive.

  \vspace{2mm}
  \textit{Hint:} Using the strategy from class where $h=k$, consider a decomposition of
  the request sequence $\sigma$ into phases that are maximal subsequences with $k$
  distinct pages requested.

\end{document}

https://www.ibr.cs.tu-bs.de/courses/ss18/oa/exercise_sheets/ex4.pdf
in case we ever want some harder problems